\chapter{測量、么正演化與量子動力學}

\lectureoverview{我們先整理已經用到的測量原理：可觀測量是厄米算符，其本徵值是可能的測量結果，可區分的態彼此正交，而玻恩定則給出機率。接下來要問的是動力學。可逆性要求所有初始態都遵循同一個線性演化，且演化保持內積不變；因此演化算符是么正算符，其無窮小生成元則是厄米算符。這個生成元就是哈密頓量，其作用給出含時Schr\"odinger方程。接著我們推導期望值公式、描述平均值演化的對易子方程、此方程與經典泊松括號的關係，以及平均能量守恆。}

暫且先把自旋系統放在一邊，等一般結構建立之後再回來討論。以下原理並非一組彼此獨立且不可再精簡的原理。有些原理彼此補充或部分蘊含；分別陳述，則能更清楚地呈現各自的物理意涵。

\section{測量可能得到什麼}

\emph{可觀測量}就是可以測量的物理量。沿用至今的這個術語稍嫌隆重；\emph{可測量量}其實也同樣適用。每個可觀測量都對應一個厄米算符。暫且以\(L\)表示這個算符：
\begin{equation}
L^\dagger=L.
\end{equation}
並非每個線性算符都是厄米算符。反過來說，能否用現有材料和儀器測量某個特定的厄米算符，可能是個棘手的實驗問題。不過，抽象的工作規則相當寬鬆：厄米算符代表可能的可觀測量，即使實際測量方法尚未發明也是如此。

如果
\begin{equation}
L\ket{\lambda}=\lambda\ket{\lambda},
\end{equation}
那麼，\(\lambda\)就是測量\(L\)時可能得到的數值結果。我們用同一個符號標記本徵值及其本徵向量。右矢的括號能將兩者區分開來：\(\lambda\)是數，而\(\ket{\lambda}\)是態。

為什麼要用厄米算符？最基本的測量儀器回報的是實數。如果儀器看似回報一個複數，實際上它回報的是兩個獨立的實數。因此，可觀測量的本徵值必須是實數。此外，對應不同確定結果的態必須在物理上可區分，因此彼此正交。在此處採用的有限維情形中，若線性算符具有完備的正交歸一本徵基，且所有本徵值都是實數，則該算符是厄米算符。因此，物理陳述與厄米結構彼此吻合。

\subsection{可區分性就是正交性}

\emph{物理上可區分}有明確的操作意義。假設一個自旋被製備為向上或向下。沿\(z\)軸測量，就能毫無歧義地分辨兩種可能：向上得到\(+1\)，向下得到\(-1\)。同樣地，沿\(x\)軸測量可以區分左與右。

現在假設自旋被製備為向上或向右。沿\(z\)軸測量時，向上態必定得到\(+1\)，但向右態也有非零機率得到\(+1\)。沿\(x\)軸測量，只會把這種不對稱互換。將儀器傾斜\(45^\circ\)，也無法消除兩種結果分佈的重疊。無論儀器朝向如何，都無法得到一個總能區分這兩種製備的結果。

\begin{definition}
若某個實驗或一組實驗能毫無歧義地判定所提供的是哪一種製備，這兩個態在物理上便可區分。
\end{definition}

量子力學以幾何方式表示完全可區分性：
\begin{equation}
\text{distinguishable states}
\quad\Longrightarrow\quad
\text{orthogonal vectors}.
\end{equation}

\begin{classroomqa}
\audiencequestion{換一條測量軸，能否區分向上態和向右態？}
\lecturerresponse{不能。沿\(z\)或\(x\)軸測量，有時可以判別是哪一種製備，但至少有一種結果仍與兩種製備都相容。沿中間方向測量，機率分佈也仍有重疊。問題不在於儀器選得不好，而在於這兩個態向量不正交。}
\end{classroomqa}

\section{玻恩定則}

本徵值告訴我們可能發生什麼，卻沒有告訴我們各結果出現的機率。機率規則提供了預測所需的資訊。如果系統被製備在歸一化態\(\ket{A}\)，而歸一化本徵向量\(\ket{\lambda}\)屬於所測量的可觀測量\(L\)，那麼
\begin{equation}
P(\lambda)=\left|\braket{\lambda}{A}\right|^2.
\label{eq:l4-born-rule}
\end{equation}
重疊\(\braket{\lambda}{A}\)稱為\emph{機率振幅}。它通常是複數，因此本身不能是機率；其絕對值平方則是實數且非負。

\begin{figure}[htbp]
\centering
\includegraphics[width=0.86\textwidth]{lecture_04_figure_01.png}
\caption{玻恩定則與測量原理並列：機率等於重疊的絕對值平方。}
\lectureframe{00:12:40}
\end{figure}

選取\(L\)的一組正交歸一本徵基，並將態展開為
\begin{equation}
\ket{A}=\sum_i \alpha_i\ket{\lambda_i},
\qquad
\alpha_i=\braket{\lambda_i}{A}.
\end{equation}
於是玻恩定則成為
\begin{equation}
P(\lambda_i)=|\alpha_i|^2.
\end{equation}
歸一化條件則給出
\begin{equation}
\sum_i P(\lambda_i)=\sum_i|\alpha_i|^2=1.
\end{equation}
正定性和歸一化條件本身無法推導出玻恩定則，但它們說明絕對值平方為何如此自然地符合機率論。歸根究柢，這條定則仍須由實驗檢驗。

\begin{classroomqa}
\audiencequestion{為了讓所有機率加總為一，是否也必須取絕對值平方？}
\lecturerresponse{歸一化態以正交歸一基展開時，各振幅的平方和確實為一。因此，正定性與歸一化條件都與玻恩定則完全吻合。在此處，哪個陳述在邏輯上優先並不如兩者彼此一致來得重要；機率規則由實驗確立。}
\end{classroomqa}

\begin{classroomqa}
\audiencequestion{本徵值是數，為什麼要把它寫在右矢裡？}
\lecturerresponse{記號\(\ket{\lambda}\)表示對應於數值本徵值\(\lambda\)的本徵向量。若要更嚴謹，可以另外引入一個標籤，但指標一多，計算很快就會變得難以理解。右矢符號能讓數與向量有所區別。}
\end{classroomqa}

\begin{classroomqa}
\audiencequestion{是哪條原理讓多種可能結果中的一種成為實際結果？}
\lecturerresponse{玻恩定則給出各結果的機率。如果製備的態已是\(L\)的本徵向量，對應的結果便以機率一出現。若態是疊加態，則會依照方程~\eqref{eq:l4-born-rule}所列的機率得到其中一種結果。相關的態更新或坍縮規則是另一個測量問題；等討論糾纏之後，我們會再回來探討。}
\end{classroomqa}

\subsection{簡併的測量結果}

一個本徵值可能對應不只一個線性獨立的本徵向量。假設
\begin{align}
L\ket{\lambda,1}&=\lambda\ket{\lambda,1},\\
L\ket{\lambda,2}&=\lambda\ket{\lambda,2}.
\end{align}
它們張成的空間內，任意線性組合也都是本徵值為\(\lambda\)的本徵向量：
\begin{equation}
L\bigl(\alpha\ket{\lambda,1}+\beta\ket{\lambda,2}\bigr)
=
\lambda\bigl(\alpha\ket{\lambda,1}+\beta\ket{\lambda,2}\bigr).
\end{equation}
這項自由度讓我們能在簡併本徵空間內選擇一組正交歸一基。若標籤\(a\)遍歷這組基，數值結果\(\lambda\)的機率就是各個互斥正交通道的機率總和：
\begin{equation}
P(\lambda)=\sum_a\left|\braket{\lambda,a}{A}\right|^2.
\label{eq:l4-degenerate-born}
\end{equation}
等價地，若\(P_\lambda\)是投影到整個本徵空間的投影算符，則\(P(\lambda)=\bra{A}P_\lambda\ket{A}\)。投影算符的記法較簡潔；方程~\eqref{eq:l4-degenerate-born}則直接呈現機率如何相加。

\begin{classroomqa}
\audiencequestion{如果製備態與多個本徵向量都有非零重疊，或者多個本徵向量具有相同本徵值，該怎麼辦？}
\lecturerresponse{本徵值不同時，每個重疊都給出各自測量結果的機率。若多個正交本徵向量具有相同本徵值，就將各重疊的絕對值平方相加。如果一開始選取的簡併子空間基底並不正交，先將它們換成正交歸一基；相加所得的機率與所選基底無關。}
\end{classroomqa}

\begin{classroomqa}
\audiencequestion{若完全不知道自旋是如何製備的，能否得到明確答案？}
\lecturerresponse{只拿一個任意樣本做不到。若對製備方式沒有任何限制，就不存在一種測量能毫無歧義地辨認未知量子態。製備資訊會限定測量能夠區分哪些可能性。}
\end{classroomqa}

\section{新問題：態如何改變}

前面的原理描述了態與測量。動力學要問的是，從一個時刻到另一個時刻之間會發生什麼。經典類比很有幫助。兩態硬幣可能每次都保持不變，也可能依序交替出現正面、反面、正面、反面。六面骰子可以有更多更新規則，而連續狀態空間則有無限多種可能。動力學定律會把目前的態映射到下一個態。

令人滿意的經典定律應具決定性且可逆。不同的態不會合併成同一個態，否則僅憑現在的狀態就無法確定過去。這種資訊的保存有時稱為\emph{負第一定律}：它先於較熟悉的守恆定律。

量子理論從態向量\(\ket{\psi(t)}\)出發，並明確將它視為時間的函數。假設存在一個線性的時間演化算符\(U(t)\)，使得
\begin{equation}
\ket{\psi(t)}=U(t)\ket{\psi(0)}.
\label{eq:l4-time-development}
\end{equation}
同一個\(U(t)\)作用於系統的每一個初始態。它不會根據輸入向量而重新選定，且其作用是線性的。更一般的非線性或依賴態的定律並非不能設想，但它們很快就會與量子力學的機率、因果關係及複合系統結構相衝突。

\subsection{可逆性保留內積}

令另一個初始態依照相同定律演化：
\begin{equation}
\ket{\phi(t)}=U(t)\ket{\phi(0)}.
\end{equation}
若兩個製備在初始時刻可以區分，演化後也必須仍可區分。因此
\begin{equation}
\braket{\psi(0)}{\phi(0)}=0
\quad\Longrightarrow\quad
\braket{\psi(t)}{\phi(t)}=0.
\end{equation}

\begin{figure}[htbp]
\centering
\includegraphics[width=0.86\textwidth]{lecture_04_figure_02.png}
\caption{兩個態受同一個\(U(t)\)演化；初始的正交性在演化後仍然保留。}
\lectureframe{00:34:43}
\end{figure}

只要演化能保留態的歸一化和正交性，就足以保留所有內積。為了看清箇中機制，將兩個任意向量以一組正交歸一基展開。演化後，基底向量仍保持正交歸一，因此各係數在每個時刻重組出的重疊都相同：
\begin{equation}
\braket{\psi(t)}{\phi(t)}=\braket{\psi(0)}{\phi(0)}.
\label{eq:l4-overlap-conservation}
\end{equation}
因此，演化不只保留完全可區分性，也保留兩個態之間完整的相似程度。

對式~\eqref{eq:l4-time-development}取厄米共軛，可得對應的左矢：
\begin{equation}
\bra{\psi(t)}=\bra{\psi(0)}U^\dagger(t).
\end{equation}
將左矢與右矢的演化式代入式~\eqref{eq:l4-overlap-conservation}，得到
\begin{equation}
\bra{\psi(0)}U^\dagger(t)U(t)\ket{\phi(0)}
=
\braket{\psi(0)}{\phi(0)}
\end{equation}
此式對每一對初始向量都成立。若某個算符在任意兩個向量間的矩陣元都與恆等算符相同，它就是恆等算符。因此
\begin{equation}
U^\dagger(t)U(t)=I.
\label{eq:l4-unitarity}
\end{equation}
滿足式~\eqref{eq:l4-unitarity}的算符稱為\emph{么正算符}。在有限維情形下，
\begin{equation}
U^\dagger(t)=U^{-1}(t),
\end{equation}
所以演化也可以反向進行。

\begin{classroomqa}
\audiencequestion{此處的大寫與小寫\(\Phi\)代表不同種類的態嗎？}
\lecturerresponse{不是，那只是書寫不一致。我們一律使用\(\ket{\phi}\)；在數學上，態由其標籤和取值區分，而非偶然的大小寫。}
\end{classroomqa}

\begin{classroomqa}
\audiencequestion{可以把\(U^\dagger U=I\)解讀為其中一個因子是另一個因子的反算符嗎？}
\lecturerresponse{可以。么正算符的反算符是\(U^{-1}=U^\dagger\)。這就是可逆且保留內積的演化在代數上的表述。}
\end{classroomqa}

\section{無窮小時間與哈密頓量}

經過時間為零時，態不會改變，因此\(U(0)=I\)。在極短的時間間隔\(\epsilon\)內，\(U(\epsilon)\)必須接近恆等算符。暫時採用令\(\hbar=1\)的單位制，寫成
\begin{equation}
U(\epsilon)=I-i\epsilon H+O(\epsilon^2).
\label{eq:l4-small-u}
\end{equation}
負號與因子\(i\)起初只是記號約定；尚未知道\(H\)時，它們可以併入該算符的定義。施加么正條件後，這種寫法的好處就會顯現。取厄米共軛可得
\begin{equation}
U^\dagger(\epsilon)=I+i\epsilon H^\dagger+O(\epsilon^2).
\end{equation}
於是
\begin{align}
I
&=U^\dagger(\epsilon)U(\epsilon)\\
&=\bigl(I+i\epsilon H^\dagger\bigr)
  \bigl(I-i\epsilon H\bigr)+O(\epsilon^2)\\
&=I+i\epsilon(H^\dagger-H)+O(\epsilon^2).
\end{align}
一階項必須為零：
\begin{equation}
H^\dagger=H.
\label{eq:l4-h-hermitian}
\end{equation}

\begin{figure}[htbp]
\centering
\includegraphics[width=0.86\textwidth]{lecture_04_figure_03.png}
\caption{一階么正性計算要求無窮小生成元為厄米算符。}
\lectureframe{00:43:38}
\end{figure}

字母\(H\)代表的不是「厄米」，而是\emph{哈密頓量}。由於它是厄米算符，因此代表一個可觀測量。它被認定為能量；從更具結構性的觀點來看，它是時間平移的生成元。這項認定不只是名稱上的約定：對於經典極限已知的系統，量子\(H\)必須再現經典哈密頓量所描述的動力學。對稱性、自由度、物理判斷，最終還有實驗，共同決定哪個厄米算符是特定系統的哈密頓量。

\begin{classroomqa}
\audiencequestion{任何厄米算符都能作為\(H\)嗎？實際採用哪一個由什麼決定？}
\lecturerresponse{目前的推導容許希爾伯特空間上的任意厄米算符。物理上採用的算符則由系統的自由度建構而成。以自旋系統為例，它必須由自旋算符和外場建構。對稱性和已知的相互作用會縮小可能範圍；實驗則確定剩餘的常數，有時甚至能決定算符的形式。}
\end{classroomqa}

哈密頓量可以顯式地依賴時間。電流隨時間改變的磁場就是一個簡單例子。若外部參數不隨時間改變，則\(H\)沒有顯式的時間依賴；這就是時間平移不變性。若外部參數隨時間變化，能量便不一定守恆。

\begin{classroomqa}
\audiencequestion{在小時間間隔的展開中，\(H\)是常數嗎？還是應該寫成\(H(0)\)？}
\lecturerresponse{對於不隨時間改變的問題，\(H\)在所有時刻都是同一個算符。更一般地，從時刻\(t\)開始的局部演化步驟會包含\(H(t)\)。此處的推導先處理時間平移不變的情形，再說明顯式時間依賴如何影響能量守恆。}
\end{classroomqa}

\section{含時Schr\"odinger方程}

將式~\eqref{eq:l4-small-u}作用於態：
\begin{equation}
\ket{\psi(\epsilon)}=(I-i\epsilon H)\ket{\psi(0)}+O(\epsilon^2).
\end{equation}
整理後除以\(\epsilon\)，得到
\begin{equation}
\frac{\ket{\psi(\epsilon)}-\ket{\psi(0)}}{\epsilon}
=-iH\ket{\psi(0)}+O(\epsilon).
\label{eq:l4-difference-quotient}
\end{equation}

\begin{figure}[htbp]
\centering
\includegraphics[width=0.86\textwidth]{lecture_04_figure_04.png}
\caption{將無窮小態更新式整理成向量值的差商。}
\lectureframe{00:51:42}
\end{figure}

對向量求導並沒有什麼特別神祕之處。兩個向量可以相減，其差可以除以一個數，而取極限後便定義出另一個向量。前面將參考時刻稱為零，這個時刻其實可以是任意瞬間。因此令\(\epsilon\to0\)，便得到
\begin{equation}
\boxed{\frac{d}{dt}\ket{\psi(t)}=-iH\ket{\psi(t)}}
\qquad(\hbar=1).
\label{eq:l4-schrodinger-natural}
\end{equation}

\begin{figure}[htbp]
\centering
\includegraphics[width=0.86\textwidth]{lecture_04_figure_05.png}
\caption{黑板上的推導完成後，最後得到方框中的含時Schr\"odinger方程。}
\lectureframe{00:56:11}
\end{figure}

這就是一般形式的含時Schr\"odinger方程。熟悉的粒子波函數微分方程，是這個抽象態向量方程的一種表示方式。要使用此方程，必須知道\(H\)和初始態。

若要描述有限時間間隔，就要求解這個微分方程。當\(H\)不隨時間改變時，反覆複合無窮小演化步驟可得
\begin{equation}
U(t)=e^{-iHt}
\qquad(\hbar=1),
\end{equation}
這正是複合所有小步驟的精確表述。若哈密頓量隨時間改變，則需要引入時間排序；這個進一步的細節留待稍後討論。

\begin{classroomqa}
\audiencequestion{前面的展開用的是無窮小時間間隔。經過長時間後的態要如何決定？}
\lecturerresponse{反覆套用無窮小演化，或等價地求解微分方程。若\(H\)不隨時間改變，所有小步驟的乘積就成為\(e^{-iHt/\hbar}\)。}
\end{classroomqa}

\subsection{恢復普朗克常數}

前面的方程採用\(\hbar=1\)的單位制。量綱分析可以指出\(\hbar\)應出現的位置。時間導數的量綱是時間的倒數，而\(H\)的量綱是能量。由於\(\hbar\)的量綱是能量乘以時間，
\begin{equation}
\boxed{i\hbar\frac{d}{dt}\ket{\psi(t)}=H\ket{\psi(t)}}.
\label{eq:l4-schrodinger}
\end{equation}
等價地，無窮小演化算符為\(U(\epsilon)=I-i\epsilon H/\hbar+O(\epsilon^2)\)，而不隨時間改變的有限時間演化算符為\(U(t)=e^{-iHt/\hbar}\)。

\begin{classroomqa}
\audiencequestion{普朗克常數確切出現在哪裡？}
\lecturerresponse{普朗克常數以以下兩種等價形式之一出現：
\begin{align*}
i\hbar\,\partial_t\ket{\psi}&=H\ket{\psi},\\
\partial_t\ket{\psi}&=-\frac{i}{\hbar}H\ket{\psi}.
\end{align*}
這兩式的量綱都相符，因為\(\hbar\)的單位是能量乘以時間。}
\end{classroomqa}

態依照方程~\eqref{eq:l4-schrodinger}確定性地演化；測量結果仍具有統計性，因為玻恩定則會從不斷變化的態中抽取測量結果。

\begin{classroomqa}
\audiencequestion{這項推導是否仰賴可逆性？若測量會擾動態，推導又如何成立？}
\lecturerresponse{推導仰賴線性演化和內積守恆。此方程描述的是孤立系統。測量儀器與系統的交互作用強到足以使儀器成為物理系統的一部分，因此不能忽略。必須以量子力學處理系統與儀器的組合；在分析測量之前，必須先理解兩者之間的糾纏。}
\end{classroomqa}

\section{期望值就是平均值}

若要將量子動力學與經典方程相比較，就需要可觀測量的平均值。標準術語\emph{期望值}可能會讓人誤解。對於只取\(+1\)或\(-1\)的公平變數，平均值是零，雖然零從未出現過。對稱分佈甚至可能在均值附近一整段區域內都為零。某些平滑的單峰分佈中，均值、中位數和眾數會重合，但一般而言並非如此。

若\(L\)可能測得數值\(\lambda_i\)，各自的機率為\(P(\lambda_i)\)，則定義
\begin{equation}
\overline{\lambda}=\sum_i\lambda_iP(\lambda_i).
\label{eq:l4-probability-average}
\end{equation}
這是分佈的第一動差，不代表單次測量會得到這個結果。

\begin{classroomqa}
\audiencequestion{最常出現的值不才是最值得期待的嗎？}
\lecturerresponse{最常出現的值稱為眾數，加權平均則是均值。在簡單而集中的分佈中，兩者可能相同；但雙峰分佈的均值可能落在機率幾乎為零的區域。「期望」指的是數學上的均值，不一定是最可能出現的單一結果。}
\end{classroomqa}

重複試驗能將數學上的均值與實驗資料連結起來。在通常的獨立性和動差有限等假設下，大數定律描述樣本平均值依機率收斂的性質。它並不保證每一段有限長度的結果都接近均值；連續出現極長一串相同自旋結果仍有可能，只是機率極低。這也是一種統計上的微妙之處，使熱力學第二定律並非毫無例外的微觀禁令。

\begin{classroomqa}
\audiencequestion{如果公平的自旋實驗碰巧連續出現一百萬次向上，這段結果的平均值真的是向上嗎？}
\lecturerresponse{這段特定結果的樣本平均值確實是向上。理論期望值仍然是零。隨著獨立試驗次數增加，樣本平均值以高機率趨近理論均值，但機率論並不會排除所有特殊的有限序列。}
\end{classroomqa}

\subsection{算符夾心式}

將歸一化態以\(L\)的一組正交歸一本徵基展開：
\begin{equation}
\ket{A}=\sum_i\alpha_i\ket{\lambda_i},
\qquad
\bra{A}=\sum_j\alpha_j^*\bra{\lambda_j}.
\end{equation}
根據玻恩定則和式~\eqref{eq:l4-probability-average}，可得
\begin{equation}
\overline{\lambda}=\sum_i\lambda_i|\alpha_i|^2.
\end{equation}
現在計算
\begin{align}
\bra{A}L\ket{A}
&=\sum_{i,j}\alpha_j^*\alpha_i
  \bra{\lambda_j}L\ket{\lambda_i}\\
&=\sum_{i,j}\alpha_j^*\alpha_i\lambda_i
  \braket{\lambda_j}{\lambda_i}\\
&=\sum_i\lambda_i|\alpha_i|^2.
\end{align}
由於\(\braket{\lambda_j}{\lambda_i}=\delta_{ji}\)，雙重總和中的項便約化。因此
\begin{equation}
\boxed{\expect{L}_A=\bra{A}L\ket{A}}.
\label{eq:l4-expectation-sandwich}
\end{equation}

\begin{figure}[htbp]
\centering
\includegraphics[width=0.86\textwidth]{lecture_04_figure_06.png}
\caption{本徵基展開將機率平均值化為算符夾心式\(\bra{A}L\ket{A}\)。}
\lectureframe{01:11:50}
\end{figure}

平均值有時用尖括號表示，有時則用上橫線表示。當態隨時間演化時，
\begin{equation}
\expect{L}_t=\bra{\psi(t)}L\ket{\psi(t)}.
\end{equation}
下標只是提醒我們，態及其平均值都依賴時間。在此處的Schr\"odinger繪景討論中，\(L\)本身是固定的，除非特別指出它有顯式時間依賴。

\begin{classroomqa}
\audiencequestion{\(t\)應該寫在括號外嗎？上橫線表示樣本平均值嗎？}
\lecturerresponse{兩種記法都可以加上標記以呈現時間依賴，但標準的量子力學記法是\(\expect{L}_t\)，或簡寫為\(\expect{L}\)，並由\(\ket{\psi(t)}\)表示時間。此處的上橫線代表理論上的機率平均值，也就是第一動差\(\sum_i\lambda_iP(\lambda_i)\)，並非特指有限筆資料的樣本平均值。}
\end{classroomqa}

\section{平均值如何演化}

用乘積微分法則對夾心式求導：
\begin{equation}
\frac{d}{dt}\expect{L}
=
\left(\frac{d}{dt}\bra{\psi}\right)L\ket{\psi}
+
\bra{\psi}L\left(\frac{d}{dt}\ket{\psi}\right).
\label{eq:l4-product-derivative}
\end{equation}
在\(\hbar=1\)單位制下，右矢方程及其厄米共軛為
\begin{align}
\frac{d}{dt}\ket{\psi}&=-iH\ket{\psi},\\
\frac{d}{dt}\bra{\psi}&=+i\bra{\psi}H.
\end{align}
複共軛會改變\(i\)的符號，而且\(H^\dagger=H\)。代入式~\eqref{eq:l4-product-derivative}可得
\begin{align}
\frac{d}{dt}\expect{L}
&=i\bra{\psi}HL\ket{\psi}
-i\bra{\psi}LH\ket{\psi}\\
&=i\bra{\psi}(HL-LH)\ket{\psi}.
\end{align}

算符不一定彼此對易，順序就如矩陣乘法一樣重要。定義
\begin{equation}
[H,L]=HL-LH.
\end{equation}
於是
\begin{equation}
\boxed{\frac{d}{dt}\expect{L}=i\expect{[H,L]}}
\qquad(\hbar=1).
\label{eq:l4-average-evolution-natural}
\end{equation}

\begin{figure}[htbp]
\centering
\includegraphics[width=0.86\textwidth]{lecture_04_figure_07.png}
\caption{平均值的時間導數由對易子的期望值決定。}
\lectureframe{01:22:22}
\end{figure}

恢復量綱後，
\begin{equation}
\boxed{\frac{d}{dt}\expect{L}
=\frac{i}{\hbar}\expect{[H,L]}
=-\frac{i}{\hbar}\expect{[L,H]}}.
\label{eq:l4-average-evolution}
\end{equation}
若\(L\)本身顯式依賴時間，一般方程還會多出一項\(\expect{\partial L/\partial t}\)。此處的\(L\)是固定的。

\begin{classroomqa}
\audiencequestion{對易子方程中的符號確實是正\(i\)嗎？}
\lecturerresponse{是的，根據\([H,L]=HL-LH\)，結果是\(+(i/\hbar)\expect{[H,L]}\)。反轉對易子的順序會使符號反轉，因此也可寫成等價形式\(-(i/\hbar)\expect{[L,H]}\)。}
\end{classroomqa}

\begin{classroomqa}
\audiencequestion{左矢方程是由複共軛還是厄米共軛得到的？為什麼導數仍留在左邊？}
\lecturerresponse{這裡取的是厄米共軛。右矢的導數仍是右矢，而該導數的對偶就是對應左矢的導數。從差商也能直接看出這一點：右矢之差仍是右矢，左矢之差也仍是左矢。時間在此是參數，不是必須反轉運算順序的算符。}
\end{classroomqa}

\section{經典的映照}

在經典哈密頓力學中，不顯式依賴時間的相空間函數\(L(q,p)\)滿足
\begin{equation}
\frac{dL}{dt}=\{L,H\}_{\mathrm{PB}},
\label{eq:l4-poisson-evolution}
\end{equation}
其中\(\{\cdot,\cdot\}_{\mathrm{PB}}\)是泊松括號。將此式與式~\eqref{eq:l4-average-evolution}相比。兩種括號都具有反對稱性：
\begin{align}
\{L,H\}_{\mathrm{PB}}&=-\{H,L\}_{\mathrm{PB}},\\
[L,H]&=-[H,L].
\end{align}
兩者也都像導數一樣作用於乘積：
\begin{align}
\{LM,H\}_{\mathrm{PB}}
&=\{L,H\}_{\mathrm{PB}}M+L\{M,H\}_{\mathrm{PB}},\\
[LM,H]
&=L[M,H]+[L,H]M.
\label{eq:l4-commutator-product}
\end{align}
在\(LHM\)上加一項並減去同一項，即可得到對易子恆等式：
\begin{align}
LMH-HLM
&=(LMH-LHM)\\
&\quad +(LHM-HLM)\\
&=L(MH-HM)\\
&\quad +(LH-HL)M.
\end{align}

這種平行關係提示我們，兩者之間存在以下對應：
\begin{equation}
[A,B]
\quad\longleftrightarrow\quad
i\hbar\{A,B\}_{\mathrm{PB}},
\label{eq:l4-dirac-correspondence}
\end{equation}
等價地，也可以寫成
\begin{equation}
\{A,B\}_{\mathrm{PB}}
\quad\longleftrightarrow\quad
-\frac{i}{\hbar}[A,B].
\end{equation}

\begin{figure}[htbp]
\centering
\includegraphics[width=0.86\textwidth]{lecture_04_figure_08.png}
\caption{經典泊松括號方程與量子對易子的對應關係，並已恢復\(\hbar\)。}
\lectureframe{01:32:32}
\end{figure}

狄拉克注意到這項關聯，並以此作為理解量子力學的途徑。不過，這應視為一種對應關係，而不是普遍適用的代換規則：算符排序可能產生歧義，而經典極限必須由量子動力學推導而來。從概念上說，量子力學是基礎；對於適當的態和尺度，經典力學是量子力學的近似。對易子在理論上先於泊松括號，儘管歷史上的類比有助於我們發現兩者的關係。

\begin{classroomqa}
\audiencequestion{為什麼把\(H\)稱為能量？為什麼在提出對稱性原理之前，就推論能量守恆？}
\lecturerresponse{在此階段，\(H\)嚴格來說是時間平移的厄米生成元。將它認定為能量，是一項物理上的識別；對於已知經典極限的系統，必須檢驗這項識別是否成立。對於處於適當態的大型或重型系統，量子平均值必須遵循經典方程，而量子\(H\)必須與經典哈密頓量相連結。以下的守恆陳述假設這個生成元沒有顯式時間依賴。}
\end{classroomqa}

\section{平均能量守恆}

在哈密頓量不隨時間改變的條件下，將\(L=H\)代入方程~\eqref{eq:l4-average-evolution}：
\begin{equation}
\frac{d}{dt}\expect{H}
=\frac{i}{\hbar}\expect{[H,H]}.
\end{equation}
任何算符都與自身對易：
\begin{equation}
[H,H]=HH-HH=0,
\end{equation}
因此
\begin{equation}
\boxed{\frac{d}{dt}\expect{H}=0}.
\label{eq:l4-energy-conservation}
\end{equation}

\begin{figure}[htbp]
\centering
\includegraphics[width=0.86\textwidth]{lecture_04_figure_09.png}
\caption{哈密頓量與自身對易，因此平均能量立即可知是守恆的。}
\lectureframe{01:39:30}
\end{figure}

這證明了一個刻意限縮的結論：平均能量保持不變。相同的動力學也將導出關於整個能量分佈的更強結論。如果\(H\)明確依賴時間，額外的顯式時間導數就不一定為零；受驅動的系統可以和任何使其參數改變的因素交換能量。

還有兩個問題懸而未決：沿任意軸測量自旋的機率仍待計算，狀態更新問題也尚未推導出來。在兩次干預之間，孤立系統的狀態會作么正演化。測量時，系統與測量儀器相互作用並形成糾纏。

\begin{classroomqa}
\audiencequestion{之前說好要計算沿任意軸測量自旋的機率，後來怎麼了？}
\lecturerresponse{一般動力學佔去了可用的時間。自旋計算仍是下一個具體的應用，待這套結構框架建立後就會接著處理。}
\end{classroomqa}

\begin{classroomqa}
\audiencequestion{沒有坍縮公設，這些公設怎麼可能完備？}
\lecturerresponse{目前所說的坍縮只是暫時的說法。測量會使系統與測量儀器形成糾纏。在複合系統與糾纏的概念尚未建立之前，無法如實建模測量過程。下一步就是建立這個概念。}
\end{classroomqa}
